\begin{ThreePartTable}

\begin{TableNotes}
    \item[*] 0x4xxx\_xxxx 开头的地址可配置为 cache-able 或 noncache-able，具体由 CPU PMA 控制。0x8xxx\_xxxx 开头的地址为 CPU 直接访问。
    \item[*] 关于 HP CPU 和 LP CPU 可以访问的目标，请参考图 \ref{fig:sysmem-address-mapping-structure}。
\end{TableNotes}
\begin{longtable}[c]{|p{3.5cm}|c|c|r|c| }
    \caption{地址映射}
    \label{tab:sysmem-memory}
       \\\hline
        \rowcolor{lightgray}
       & \multicolumn{2}{c|}{\textbf{边界地址}} & & \\\cline{2-3}

       \rowcolor{lightgray}
       \multirow{-2}*{\textbf{总线类型}} & \textbf{低位地址} & \textbf{高位地址} & \multirow{-2}*{\textbf{容量}} & \multirow{-2}*{\textbf{目标\tnote{*}}} \\\hline
    \endfirsthead

       \multicolumn{5}{c}{\bfseries \tablename\ \thetable{} -- 接上页}\\\hline

       \rowcolor{lightgray}
       & \multicolumn{2}{c|}{\textbf{边界地址}} & & \\\cline{2-3}

       \rowcolor{lightgray}
       \multirow{-2}*{\textbf{总线类型}} & \textbf{低位地址} & \textbf{高位地址} & \multirow{-2}*{\textbf{容量}} & \multirow{-2}*{\textbf{目标\tnote{*}}} \\\hline
    \endhead

        \multicolumn{5}{r}{\textbf{见下页}} \\
    \endfoot
        \insertTableNotes
    \endlastfoot

        \rowcolor{gray!25}  & 0x0000\_0000 & 0x300F\_FFFF &        & 保留      \\\hline
        {数据总线/指令总线} & 0x3010\_0000 & 0x3010\_1FFF & 8 KB  & HP SPM  \\\hline
        \hline
        \rowcolor{gray!25}  & 0x3010\_2000 & 0x3FEF\_FFFF &         & 保留      \\\hline
        {数据总线/指令总线} & 0x3FF0\_0000 & 0x3FF1\_FFFF & 128 KB  & HP CPU 外设 \\\hline
        \hline
        \rowcolor{gray!25}  & 0x3FF2\_0000 & 0x3FFF\_FFFF &         & 保留      \\\hline
        {数据总线/指令总线} & 0x4000\_0000 & 0x43FF\_FFFF & 64 MB & 外部 flash \\\hline
        \rowcolor{gray!25}  & 0x4400\_0000 & 0x47FF\_FFFF &         & 保留      \\\hline
        {数据总线/指令总线} & 0x4800\_0000 & 0x4BFF\_FFFF & 64 MB  & 外部 RAM
        \\\hline
        \rowcolor{gray!25}  & 0x4C00\_000 & 0x4FBF\_FFFF &         & 保留      \\\hline
        {数据总线/指令总线} & 0x4FC0\_0000 & 0x4FC1\_FFFF & 128 KB  & HP ROM
        \\\hline
        \rowcolor{gray!25}  & 0x4FC2\_0000 & 0x4FEF\_FFFF &         & 保留      \\\hline
        {数据总线/指令总线} & 0x4FF0\_0000 & 0x4FFB\_FFFF & 768 KB  & HP L2MEM \\\hline
        \rowcolor{gray!25}  & 0x4FFC\_0000 & 0x4FFF\_FFFF &         & 保留      \\\hline
        {数据总线/指令总线} & 0x5000\_0000 & 0x500F\_FFFF & 1 MB  & HP 外设 \\\hline
        {数据总线/指令总线} & 0x5010\_0000 & 0x5010\_3FFF & 16 KB  & LP ROM \\\hline
        \rowcolor{gray!25}  & 0x5010\_4000 & 0x5010\_7FFF &         & 保留      \\\hline
        {数据总线/指令总线} & 0x5010\_8000 & 0x5010\_FFFF & 32 KB  & LP SRAM \\\hline
        {数据总线/指令总线} & 0x5011\_0000 & 0x5012\_FFFF & 128 KB  & LP 外设 \\\hline
        \rowcolor{gray!25}  & 0x5013\_0000 & 0x7FFF\_FFFF &         & 保留      \\\hline
        {数据总线/指令总线} & 0x8000\_0000 & 0x83FF\_FFFF & 64 MB  & 外部 flash（CPU 直接访问） \\\hline
        \rowcolor{gray!25}  & 0x8400\_0000 & 0x87FF\_FFFF &         & 保留      \\\hline
        {数据总线/指令总线} & 0x8800\_0000 & 0x8BFF\_FFFF & 64 MB  & 外部 RAM（CPU 直接访问） \\\hline
        \rowcolor{gray!25}  & 0x8C00\_000 & 0x8FBF\_FFFF &         & 保留      \\\hline
        {数据总线/指令总线} & 0x8FC0\_0000 & 0x8FC1\_FFFF & 128 KB  & HP ROM（CPU 直接访问） \\\hline
        \rowcolor{gray!25}  & 0x8FC2\_0000 & 0x8FEF\_FFFF &         & 保留      \\\hline
        {数据总线/指令总线} & 0x8FF0\_0000 & 0x8FFB\_FFFF & 768 KB  & HP L2MEM（CPU 直接访问）\\\hline
        \rowcolor{gray!25}  & 0x8FFC\_0000 & 0xFFFF\_FFFF &         & 保留      \\\hline
\end{longtable}
\end{ThreePartTable}